The equation of state (EoS) of ultra-dense matter in neutron stars remains a
major unsolved problem in fundamental physics. Theoretical models predicting a
``soft'' EoS dictate that neutron stars must collapse into black holes before
reaching a mass of $2.0\,\mathrm{M}_\odot$. Consequently, a robust mass
measurement of a neutron star exceeding $2.0\,\mathrm{M}_\odot$ would
definitively rule out soft EoS models.

Radio observations of the millisecond pulsar PSR J0740+6620, located in a
binary system with a white dwarf companion, offer a unique testbed for this
prediction. According to General Relativity, as the pulsar passes through
superior conjunction, when it is directly behind its companion, the
companion's gravitational field curves spacetime, delaying the propagation of
radio pulses to Earth. This phenomenon is known as the Shapiro time delay.
Assuming an edge-on orbital inclination ($i \approx 90^\circ$) and a perfectly
circular orbit (eccentricity $e = 0$), the Shapiro delay
$\Delta t_{\mathrm{S}}$ is approximated by
\[
\Delta t_{\mathrm{S}} \approx
  \left(\frac{2GM_{\mathrm{WD}}}{c^3}\right)\left[-\ln\left(1-\cos\phi\right)\right] + b
\]
where $M_{\mathrm{WD}}$ is the white dwarf mass, $\phi$ is the orbital phase
($\phi = 0^\circ$ at superior conjunction), and $b$ is a baseline timing
offset.

The binary system has an orbital period of $P_b = 4.767$ days. The pulsar has
a projected semi-major axis of $x = 3.978$ seconds, which corresponds to an
orbital radius $a_p = cx$. Assume the uncertainties on $P_b$ and $x$ are
negligible.
% sic: "seconds" as printed (the solution file says light-seconds)

The following table presents 15 observations of pulsar Shapiro delay at
various orbital phases ($\phi$) near superior conjunction. A constant
measurement uncertainty of $\sigma_t = \pm 1.5\,\mu\mathrm{s}$ is assumed for
all observations.

\begin{center}
\begin{tabular}{ccc}
\toprule
ID & Orbital Phase $\phi$ ($^\circ$) & Shapiro Delay $\Delta t_{\mathrm{S}}$ ($\mu$s) \\
\midrule
1  & $2.0$   & 18.5 \\
2  & $-3.0$  & 17.5 \\
3  & $4.0$   & 16.3 \\
4  & $-5.0$  & 13.9 \\
5  & $6.0$   & 12.8 \\
6  & $8.0$   & 12.5 \\
7  & $10.0$  & 10.1 \\
8  & $-12.0$ & 10.3 \\
9  & $15.0$  & 9.3 \\
10 & $20.0$  & 6.5 \\
11 & $-22.0$ & 5.8 \\
12 & $25.0$  & 6.6 \\
13 & $30.0$  & 5.0 \\
14 & $-35.0$ & 4.7 \\
15 & $40.0$  & 3.9 \\
\bottomrule
\end{tabular}
\end{center}

\medskip\noindent\textbf{Part 1 (26 points): Estimation of White Dwarf Mass}

\begin{description}
\item[(a)] (4 points) Define appropriate linearized variables ($X$ and $Y$),
  and tabulate their values for 15 data points.
\item[(b)] (18 points) Plot a scatter diagram of the 15 data points including
  their available error bars. Overlay the best-fit line on your graph (mark
  your graph as ``D01.1b''). Determine the slope $a$, the intercept $b$, and
  their respective uncertainties $(\sigma_a, \sigma_b)$.
\item[(c)] (4 points) Determine the mass of the white dwarf
  $(M_{\mathrm{WD}})$ and its absolute uncertainty $(\Delta M_{\mathrm{WD}})$
  in units of $\mathrm{M}_\odot$.
\end{description}

\medskip\noindent\textbf{Part 2 (13 points): Binary Dynamics and Error Propagation}

\begin{description}
\item[(a)] (6 points) The mass function of the binary system is defined as
  \[
  f_M = \frac{M_{\mathrm{WD}}^3}{\left(M_{\mathrm{PSR}} + M_{\mathrm{WD}}\right)^2}
  \]
  where $M_{\mathrm{PSR}}$ is the mass of the pulsar. Derive the expression of
  $f_M$ in terms of the observables $a_p$ and $P_b$, and hence compute its
  value in units of $\mathrm{M}_\odot$.
\item[(b)] (5 points) Determine the mass of the pulsar $(M_{\mathrm{PSR}})$
  and derive its uncertainty $(\Delta M_{\mathrm{PSR}})$.
\item[(c)] (2 points) Determine how many standard deviations $(\sigma)$ the
  measured pulsar mass lies above the upper mass limit predicted by the
  ``soft EoS'' model. Hence, state whether the model can be ruled out.
\end{description}

\medskip\noindent\textbf{Part 3 (11 points): Pulse Time of Arrival (TOA): Effect of Multiple Factors}

In binary pulsar timing, an observed Time of Arrival (TOA) is inherently
shifted by various physical effects relative to an idealized baseline. To
accurately determine the companion mass, astrophysicists must account for
three primary components of the timing delay:
\begin{itemize}
\item Shapiro delay $(\Delta t_{\mathrm{S}})$, caused by spacetime curvature
  around the white dwarf companion, produces a minuscule shift of
  $\sim 10\,\mu\mathrm{s}$.
\item Roemer delay $(\Delta t_{\mathrm{R}})$: The variation in pulse arrival
  time caused by the movement of the pulsar along its orbit around the
  binary's center of mass.
\item Dispersion measure delay $(\Delta t_{\mathrm{DM}})$: The propagation
  delay experienced by a signal traveling through the ionized interstellar
  medium. This delay is given by
  \[
  \Delta t_{\mathrm{DM}} = K \cdot DM \cdot \frac{1}{f^2}
  \]
  where $f$ is the frequency of the observed signal, $K$ is the dispersion
  constant, and $DM$ is the dispersion measure.
\end{itemize}

\begin{description}
\item[(a)] (3 points) Physical classification: Classify each of the following
  signal delay mechanisms as achromatic (frequency-independent) or chromatic
  (frequency-dependent).

  Roemer delay: [\ ] Achromatic \quad [\ ] Chromatic\\
  Dispersion measure delay: [\ ] Achromatic \quad [\ ] Chromatic\\
  Shapiro delay: [\ ] Achromatic \quad [\ ] Chromatic
\item[(b)] (2 points) To understand the difficulty of isolating the Shapiro
  delay, we must compare the magnitudes of these effects. Given an observing
  frequency $f = 1400\,\mathrm{MHz}$ (standard L-band), an average
  interstellar dispersion measure $DM = 15\,\mathrm{pc\,cm^{-3}}$, and the
  dispersion constant
  $K = 4.15 \times 10^3\,\mathrm{s\,MHz^2\,pc^{-1}\,cm^3}$:
  \begin{description}
  \item[(b.1)] (1 points) Calculate $\Delta t_{\mathrm{DM}}$.
  \item[(b.2)] (1 points) Arrange the three delays
    $(\Delta t_{\mathrm{R}}, \Delta t_{\mathrm{DM}}, \Delta t_{\mathrm{S}})$
    in order from largest to smallest.
  \end{description}
\item[(c)] (3 points) Below are six schematic graphs (A--F) illustrating the
  time delay as a function of orbital phase over one complete orbit. The
  vertical axis represents the delay amplitude (numerical values omitted).
  Match each of the three physical delay mechanisms listed below to the graph
  that best represents its expected shape.

\ioaafig{2026_DA_D01-f1.pdf}
% f1 = six schematic graphs A to F of time delay against orbital phase (0 to 1)

  Roemer delay [\ ?\ ]\\
  Dispersion measure delay [\ ?\ ]\\
  Shapiro delay [\ ?\ ]
\item[(d)] (3 points) Observational strategy: Variations in the interstellar
  medium cause the dispersion measure to fluctuate over time, introducing
  timing noise. Select one correct answer to each question below.
  \begin{description}
  \item[(d.1)] (1 point) Why is a monochromatic observation (observing at a
    single frequency) insufficient to constrain the extremely small
    $10\,\mu\mathrm{s}$ Shapiro delay?
    \begin{description}
    \item[A.] The telescope's receiver cannot process data faster than
      $10\,\mu\mathrm{s}$.
    \item[B.] Single frequencies are completely absorbed by the interstellar
      plasma.
    \item[C.] The system contains multiple sources of time delay, so the
      different physical effects cannot be distinguished from a single
      observation.
    \item[D.] Monochromatic signals do not experience gravitational time
      dilation.
    \end{description}
  \item[(d.2)] (2 points) Which observational strategy is used to measure and
    correct for dispersion measure variations?
    \begin{description}
    \item[A.] Integrating the signal over a one-year period to average out
      the dispersion measure variations.
    \item[B.] Conducting simultaneous observations at two distinct radio
      frequencies.
    \item[C.] Observing strictly at an extremely low frequency where the
      dispersion measure can be reduced.
    \item[D.] Only observing the pulsar when it is outside superior
      conjunction.
    \end{description}
  \end{description}
\end{description}
